% 原书物理页=印刷页22--25；3.2 全节（3.2.1--3.2.4 位于第 22--23 页，3.2.5--3.2.12 位于第 23--25 页）.
\section{可解群与幂零群}\label{sec:fg-3-2}

\begin{numpara}[换位子]\label{np:fg-3-2-1}
在群 $G$ 中，两个元素 $x,y$ 的\newterm{fg-commutator}{换位子}{commutator}是指 $G$ 中的元素 $[x,y]:=x^{-1}y^{-1}xy$. \emph{对 $G$ 的两个子群 $A,B$，由 $a\in A$ 与 $b\in B$ 给出的换位子 $[a,b]$ 全体一般不是 $G$ 的子群，即使取 $A=B=G$ 时也是如此.}
\end{numpara}

\begin{definition}\label{def:fg-3-2-2}
对 $G$ 的两个子群 $A,B$，我们用 $[A,B]$ 表示 $G$ 中由换位子 $[a,b]$（$a\in A$，$b\in B$）生成的子群.
\end{definition}

几个简单的性质：
\begin{itemize}[label=--]
\item $B$ 是 $G$ 的正规子群 $\Rightarrow[A,B]\subset B$.
\item $A$ 与 $B$ 都是 $G$ 的正规子群 $\Rightarrow[A,B]$ 是 $G$ 的正规子群.
\item $A$ 与 $B$ 都是 $G$ 的特征子群 $\Rightarrow[A,B]$ 是 $G$ 的特征子群.
\end{itemize}

\begin{numpara}[换位子群]\label{np:fg-3-2-3}
子群 $DG:=[G,G]$ 称为 $G$ 的\newterm{fg-derived-subgroup}{换位子群}{derived subgroup}. 它是一个特征子群. 由其定义本身可知，$DG$ 是 $G$ 的使 $G/N$ 为 Abel 群的最小正规子群 $N$，而 $G/DG$（常记作 $G^{\mathrm{ab}}$）可以解释为 $G$ 的``最大 Abel 商''：从 $G$ 到某个 Abel 群的任何同态都唯一地通过 $G/DG$ 分解.
\end{numpara}

\begin{numpara}[导列]\label{np:fg-3-2-4}
令 $D^0G=G$，且对每个整数 $i\geqslant0$ 令 $D^{i+1}G=[D^iG,D^iG]$. 各 $D^iG$ 构成 $G$ 的一列递减的特征子群，其相继商 $D^iG/D^{i+1}G$ 都是 Abel 群；这列子群称为 $G$ 的\newterm{fg-derived-series}{导列}{derived series}. 它显然满足下面的极小性：\emph{设 $\{G_i,\ i=0,1,\dots\}$ 是一列递减的子群，满足 $G_0=G$、$G_{i+1}\triangleleft G_i$，且 $G_i/G_{i+1}$ 对每个 $i$ 都是 Abel 群. 那么，对一切 $i$ 有 $D^iG\subset G_i$.}
\end{numpara}

\begin{numpara}[可解群]\label{np:fg-3-2-5}
\emph{称 $G$ 是\newterm{fg-solvable-group}{可解群}{solvable group}，如果存在整数 $n$ 使得 $D^nG=\{1\}$. 这等价于存在一列递减的子群 $G_0=G\supset G_1\supset\cdots\supset G_n=\{1\}$，使得 $G_{i+1}\triangleleft G_i$ 且 $G_i/G_{i+1}$ 对每个 $i$ 都是 Abel 群.}
对可解群，使 $D^nG=\{1\}$ 的整数 $n$ 中最小者称为 $G$ 的\newterm{fg-solvability-class}{可解类}{solvability class}. 它也是使得存在一列递减的子群 $G_0=G\supset G_1\supset\cdots\supset G_n=\{1\}$、满足 $G_{i+1}\triangleleft G_i$ 且 $G_i/G_{i+1}$ 对每个 $i$ 都是 Abel 群的最小整数 $n$.
例如，Abel 群是可解类为 $1$ 的可解群. 二面体群 $D_n$（即 $\mathrm{O}_2(\mathbb{R})$ 中某个正 $n$ 边形的稳定子）是可解类为 $2$ 的可解群. 更一般地，由命题~\ref{prop:fg-1-1-11} 与命题~\ref{prop:fg-1-1-12} 可得
\end{numpara}

\begin{proposition}\label{prop:fg-3-2-6}
\begin{enumerate}[label=\roman*)]
\item 可解群的任何子群（相应地，任何商群）都是可解群；
\item 设 $H$ 是群 $G$ 的一个正规子群；若 $H$ 与 $G/H$ 都可解（可解类分别为 $n_1,n_2$），则 $G$ 可解（可解类 $\leqslant n_1+n_2$）.
\end{enumerate}
\end{proposition}

\emph{一点历史}：当 $n\geqslant5$ 时，交错群 $A_n$ 是单群且不是 Abel 群，因此对称群 $S_n$ 不是可解群；相反，$S_2$、$S_3$、$S_4$ 是可解群. 在 Galois 理论中人们证明，一个代数方程可用根式求解当且仅当它的 Galois 群是可解群；由此推出，当 $n\geqslant5$ 时，$n$ 次一般方程不能用根式求解，而当 $n=2$（三项式公式）、$n=3$（Cardan 公式）、$n=4$（Ferrari 公式）时则可以.

\begin{numpara}[降中心列]\label{np:fg-3-2-7}
令 $C^0G=G$，且对每个整数 $i\geqslant0$ 令 $C^{i+1}G=[G,C^iG]$. 各 $C^iG$ 构成 $G$ 的一列递减的特征子群，称为 $G$ 的\newterm{fg-lower-central-series}{降中心列}{lower central series}. \emph{对一切 $i>0$，有 $C^i(G)\supset D^i(G)$，且 $C^{i+1}(G)$ 是 $G$ 的含于 $C^i(G)$ 的最小正规子群 $N$，使得 $C^i(G)/N$ 含于 $G/N$ 的中心 $Z(G/N)$.}
\end{numpara}

\begin{numpara}[幂零群]\label{np:fg-3-2-8}
\emph{称 $G$ 是\newterm{fg-nilpotent-group}{幂零群}{nilpotent group}，如果存在整数 $n$ 使得 $C^n(G)=\{1\}$. 这些整数中最小者称为 $G$ 的\newterm{fg-nilpotency-class}{幂零类}{nilpotency class}. 群 $G$ 是幂零的，当且仅当它拥有一列子群 $G_0=G\supset G_1\supset\dots\supset G_m=\{1\}$，使得 $[G,G_i]\subset G_{i+1}$ 对每个 $i$ 成立（因为此时对一切 $i$ 有 $C^i(G)\subset G_i$）.} 这个条件等价于要求对每个 $i$ 都有 $G_i/G_{i+1}\subset Z(G/G_{i+1})$.
\end{numpara}

由此推出，幂零群的任何子群与任何商群都是幂零群. 反之，设 $H$ 是 $G$ 的含于 $G$ 的中心 $Z(G)$ 的子群，若 $G/H$ 幂零，则 $G$ 幂零.

幂零群的中心总是非平凡的. 此外，要知道一个有限群是否幂零，可以先计算它的中心，再计算 $G$ 对其中心的商群的中心，等等. 若在某一步遇到中心平凡的群，则 $G$ 不是幂零群.

\begin{example}\label{ex:fg-3-2-9}
Abel 群是幂零群. 幂零群是可解群. 但 $S_3$ 可解而不幂零（其中心是平凡的）. 同样，对任意交换域 $k$，$G=\mathrm{GL}_n(k)$ 中由上三角矩阵（对角系数任意）组成的子群 $B$ 是可解群，但当 $k\neq\mathbb{F}_2$ 时不是幂零群. 而 $B$ 中由严格上三角矩阵组成的子群 $S$ 则是幂零群.
\end{example}

\begin{proposition}\label{prop:fg-3-2-10}
\begin{enumerate}[label=\roman*)]
\item 任何 $p$-群 $G$ 都是幂零群.
\item 设 $G$ 是幂零群. $G$ 的任何真子群 $H$ 都严格包含于其正规化子 $N_G(H)$ 之中；特别地，$G$ 的任何指数为素数的子群在 $G$ 中正规.
\item 设 $G$ 是有限群，$S$ 是 $G$ 的一个 $p$-Sylow. 对 $G$ 的任何包含 $N_G(S)$ 的子群 $H$，有 $H=N_G(H)$.
\end{enumerate}
\end{proposition}

\begin{proof}
i) 根据 $G$ 通过共轭作用于其自身的类公式，$G$ 的中心 $Z$ 是非平凡的（当 $|G|>1$ 时）. 对 $|G|$ 作归纳即可得出结论. 也可以这样来证明这个结论：当 $n$ 充分大时把 $G$ 嵌入 $\mathrm{GL}_n(\mathbb{F}_p)$，然后注意到 $G$ 的各个 $p$-Sylow 都共轭于严格三角群 $S$，而后者是幂零的.
ii) 对 $G$ 的幂零类 $n$ 作归纳. $n=1$ 时结论是清楚的. 若 $n>1$，$G$ 的子群 $A=C^{n-1}(G)$ 含于 $G$ 的中心，且 $G/A$ 是幂零类为 $n-1$ 的幂零群. 于是 $N_G(H)$ 包含 $A$，若 $H$ 不包含 $A$ 则已经得证. 否则，归纳假设给出 $H/A$ 严格包含于 $N_{G/A}(H/A)$，故 $H\neq N_G(H)$.
iii) 设 $n$ 是 $N_G(H)$ 的一个元素. 那么，$S$ 与 $nSn^{-1}\subset nHn^{-1}=H$ 都是 $H$ 的 $p$-Sylow，因而在 $H$ 中相互共轭：存在 $h\in H$ 使得 $hSh^{-1}=nSn^{-1}$. 于是 $h^{-1}n\in N_G(S)\subset H$，故 $n\in H$.
\end{proof}

\begin{theorem}\label{thm:fg-3-2-11}
设 $G$ 是有限群. 下列条件等价：
\begin{enumerate}[label=\roman*)]
\item $G$ 是幂零群；
\item 对 $G$ 的任何真子群 $H$，有 $H\neq N_G(H)$；
\item 对每个素数 $p$，$G$ 拥有唯一的 $p$-Sylow；
\item $G$ 同构于其诸 Sylow 子群的直积；
\item $G$ 中阶互素的两个元素相互交换.
\end{enumerate}
\end{theorem}

\begin{proof}
i) $\Rightarrow$ ii) 已经在前面看到. 现证 ii) $\Rightarrow$ iii). 设 $S$ 是 $G$ 的一个 $p$-Sylow. 根据假设 ii) 与前一命题的 iii)，$N_G(S)=G$，也就是说 $S$ 在 $G$ 中正规. 由于我们知道所有的 $p$-Sylow 都相互共轭（定理~\ref{thm:fg-3-1-3}~ii)），故可推出 $S$ 是 $G$ 唯一的 $p$-Sylow. 再看 iii) $\Rightarrow$ iv). 唯一的 $p$-Sylow $S_p$ 必然是正规的. 由此，对每个别的素数 $\ell\neq p$，将有 $[S_p,S_\ell]\subset S_p\cap S_\ell=\{1\}$，从而得到从诸 $S_p$ 的乘积到 $G$ 的一个单同态，而由阶的理由它也是双射. 蕴含 iv) $\Rightarrow$ v) 是清楚的，其逆亦然. 蕴含 iv) $\Rightarrow$ i) 由前一命题的第~i)~点以及事实 $C^i(G_1\times G_2)=C^i(G_1)\times C^i(G_2)$ 得到.
\end{proof}

若 $G$ 是 Abel 群，它唯一的 $p$-Sylow 称为 $G$ 的\newterm{fg-p-primary-component}{$p$-准素分量}{$p$-primary component} $G^{(p)}$. 我们可以陈述：

\begin{corollary}\label{cor:fg-3-2-12}
任何有限 Abel 群 $G$ 都同构于其诸 $p$-准素分量的直积 $\prod_p G^{(p)}$，其中 $p$ 取遍 $|G|$ 的素因子全体.
\end{corollary}

此外，由 $G$ 的不变因子所给出的分解（见第 II 章，\S2）也容易导出这个结论.\footnote{译注：原文作``（voir chap. II, S2）''，疑为对另一门课程讲义的引用；本书中 $G$ 的不变因子分解位于第~\ref{chap:fg-1}~章第~\ref{sec:fg-1-3}~节.}
